% Standalone five-step Entity-Resolution pipeline figure (ITLT revision, reviewer 2, item 6).
% Build: xelatex er_pipeline.tex  → er_pipeline.pdf   (PNG: pdftoppm -r 300 -png)
\documentclass[border=6pt]{standalone}
\usepackage{fontspec}
\usepackage{polyglossia}
\setmainlanguage{ukrainian}
\setotherlanguage{english}
\setmainfont{DejaVu Serif}
\setsansfont{DejaVu Sans}
\setmonofont{DejaVu Sans Mono}
\usepackage{amsmath,amssymb}
\usepackage{tikz}
\usetikzlibrary{arrows.meta,positioning,shapes.geometric,fit,backgrounds}
\begin{document}
\begin{tikzpicture}[
  font=\small,
  node distance=5mm,
  stepbox/.style={draw, rounded corners, align=center, inner sep=5pt, minimum width=8.2cm, fill=gray!8},
  src/.style={draw, rounded corners=2pt, align=center, inner sep=4pt, minimum width=2.0cm, fill=blue!8, font=\footnotesize},
  dec/.style={draw, diamond, aspect=2.6, align=center, inner sep=1pt, fill=orange!14, font=\footnotesize},
  outbox/.style={draw, rounded corners, align=center, inner sep=6pt, minimum width=8.2cm, fill=green!12},
  side/.style={draw, rounded corners, align=center, inner sep=5pt, minimum width=3.4cm, font=\footnotesize},
  ar/.style={-{Stealth[length=2.6mm]}, thick},
  lbl/.style={font=\scriptsize, inner sep=1.5pt, fill=white}
]

% --- sources ---
\node[src] (orcid) {ORCID};
\node[src, right=4mm of orcid] (oa) {OpenAlex};
\node[src, right=4mm of oa] (sc) {Scopus};
\node[src, right=4mm of sc] (wos) {Web of Science};
\node[draw=gray!60, dashed, rounded corners, fit=(orcid)(wos), inner sep=4pt] (sources) {};
\node[font=\footnotesize\itshape, above=1mm of sources] {Сирі записи про публікації з різних джерел};

% --- steps 1-3 ---
\node[stepbox, below=7mm of sources] (s1) {\textbf{Крок 1. Канонізація сигнатури}\\[1pt]
  \footnotesize DOI $\to$ нижній регістр, зняття префіксів, валідація \;|\; назва: нормалізація NFKC, зведення регістру\\
  \footnotesize \textbf{зі збереженням кирилиці} \;|\; автори: транслітерація кирилиця\,$\to$\,латиниця};

\node[stepbox, below=of s1] (s2) {\textbf{Крок 2. Блокування} \,\footnotesize(уникнення $O(n^2)$ порівнянь)\\[1pt]
  \footnotesize блоки: точний DOI \;|\; два найдовші слова назви \;|\; «рік + слово назви»};

\node[stepbox, below=of s2] (s3) {\textbf{Крок 3. Парне логістичне зіставлення}\\[2pt]
  $s=\sigma(\boldsymbol{\beta}\cdot\boldsymbol{\phi})$\\[1pt]
  \footnotesize ознаки $\boldsymbol{\phi}$: збіг/конфлікт DOI, близькість назв, відстань Левенштейна, рік, журнал, автори, вебадреса};

% --- decision ---
\node[dec, below=8mm of s3] (d) {порівняння оцінки $s$};

% --- branches ---
\node[side, below=16mm of d, xshift=-4.9cm, fill=red!8] (drop) {\textbf{Відкинути}\\[1pt] різні праці};
\node[side, below=16mm of d, xshift=4.9cm, fill=orange!12] (hitl) {\textbf{Смуга ручної перевірки}\\[1pt] рішення ухвалює експерт};

% --- steps 4-5 ---
\node[stepbox, below=46mm of d] (s4) {\textbf{Крок 4. Кластеризація з охоронцем конфлікту DOI}\\[1pt]
  \footnotesize зв'язні компоненти за ребрами $s\ge\theta_{\text{зл}}$;\\
  \footnotesize два кластери з непорожніми неперетинними наборами DOI \textbf{не зливаються}};

\node[stepbox, below=of s4] (s5) {\textbf{Крок 5. Злиття атрибутів і фіксація походження даних}\\[1pt]
  \footnotesize голосування, зважене за надійністю джерела (Scopus, Web of Science $\succ$ OpenAlex $\succ$ ORCID)\\
  \footnotesize пріоритет полів: DOI $\succ$ Crossref $\succ$ CRIS $\succ$ OpenAlex $\succ$ ORCID \;|\; походження даних за W3C~PROV};

\node[outbox, below=of s5] (gold) {\textbf{«Золотий запис»}\\[1pt]
  \footnotesize канонічна праця з простежуваним джерелом кожного поля та оцінкою зіставлення};

% --- arrows ---
\draw[ar] (sources.south) -- (s1.north);
\draw[ar] (s1) -- (s2);
\draw[ar] (s2) -- (s3);
\draw[ar] (s3) -- (d);

% left branch: reject
\draw[ar] (d.west) -- ++(-2.2,0) node[lbl, above, pos=0.55] {$s<\theta_{\text{пер}}=0{,}70$} -| (drop.north);
% right branch: manual review band
\draw[ar] (d.east) -- ++(2.2,0) node[lbl, above, pos=0.55] {$s\in[0{,}70;\,0{,}90)$} -| (hitl.north);
% centre: auto-merge
\draw[ar] (d.south) -- node[lbl, pos=0.3, align=center] {$s\ge\theta_{\text{зл}}=0{,}90$:\\автоматичне злиття} (s4.north);
% resolved HITL decisions re-enter clustering
\draw[ar, dashed] (hitl.south) |- ([yshift=8mm]s4.north east) -- ([xshift=-8mm]s4.north east);
\node[lbl, anchor=east, font=\scriptsize\itshape] at ([xshift=-10mm, yshift=11mm]s4.north east) {рішення експерта};

\draw[ar] (s4) -- (s5);
\draw[ar] (s5) -- (gold);

\end{tikzpicture}
\end{document}
